\section[Cote 49 : un critère de platitude par injection dans des fibres plates]{Cote 49 : un critère de platitude par injection dans des fibres plates\protect\verifie[\leanblobmain]{AlgebraicGeometry/Folder49Flatness}}\label{sec:49}

La cote 49, \emph{Picard des schémas abéliens : notes manuscrites}, non datée, compte neuf pages, dont une de garde. Pour un schéma abélien $A$ projectif sur une base $S$ localement noethérienne, les feuillets~2 et~3 démontrent que le dual $A^\vee = \mathrm{Pic}^0_{A/S}$ est lisse sur $S$. Sur chaque fibre, la dimension de $A^\vee_0$ est encadrée par celle de $A_0$, par l'isogénie $\varphi_L : A \to A^\vee$ d'un côté et par le théorème de Hopf sur $H^*(A_0, \mathcal O_{A_0})$ de l'autre ; il reste à passer des fibres à la famille, ce qui demande la platitude. Le feuillet~3 note pour cela que, pour tout $x' \in A^\vee_0$, il existe $x \in A_0$ au-dessus et un homomorphisme \emph{injectif} $\mathcal O_{A^\vee_0, x'} \to \mathcal O_{A_0, x}$, puis énonce un lemme (\lot{49}{1}{3}) : $\varphi : X \to Y$ un $S$-morphisme de schémas de type fini, $F$ quasi-cohérent sur $X$ et plat sur $S$, $G$ cohérent sur $Y$, et un morphisme $G \to \varphi_*F$ dont la restriction à la fibre en $s$ est injective au-dessus de $y$ ; alors $G$ est plat sur $S$ en $y$. Les hypothèses sur $F$, $G$ et le morphisme sont ajoutées entre les lignes. La clause~c) retenue se lit \q{tout élément de $G_s$ dont l'image inverse en les $x \in \varphi_s^{-1}(y)$ s'annule est nulle}, \q{est} et \q{nulle} douteux, le second écrit au-dessus de la fin de la ligne barrée \q{$G_s \to \varphi_{s*}(F_s)$ est injectif en $y$} ; c'est la lecture de la révision du 28 novembre 2026 sur le fac-similé, passe machine non vérifiée par un lecteur. La ligne barrée se lit sans doute, et c'est elle que la lecture retient. Un argument ou une esquisse est écrit sur la même page ; la même révision en transcrit trois lignes, sous le lemme : \q{On applique ceci avec $\varphi : A \to A'$ et un pt générique de $A'$, $F = \mathcal O_A$, $G = \mathcal O_{A'}$. Cela prouve que $A'$ est plat sur $S$ cqfd.}, \q{générique} et \q{Cela} douteux. C'est l'application de la lecture, $X = A$, $Y = A^\vee$, $\varphi = \varphi_L$, $F = \mathcal O_A$, $G = \mathcal O_{A^\vee}$ ; le feuillet l'écrit donc, sans nommer $\varphi_L$.

La piste de l'onglet Findings relevait que la cible $\varphi_*F$ n'est pas cohérente sur $Y$ : le critère local de platitude usuel, où source et cible sont finies sur un même anneau local, ne s'applique pas tel quel, et la fiche se demandait s'il fallait supposer $\varphi$ fini ou quasi-fini aux points au-dessus de $y$. On démontre ici la forme locale du lemme, sans hypothèse de finitude sur $\varphi$.

\noindent\emph{Conventions.} L'énoncé est local en $y$. On pose $R = \mathcal O_{S,s}$ et $B = \mathcal O_{Y,y}$, anneaux locaux noethériens, l'homomorphisme $R \to B$ étant local ; $\mathfrak m$ est l'idéal maximal de $R$. On pose $N = G_y$, $B$-module de type fini, et, pour chaque point $x$ de $X$ au-dessus de $y$, $M_x = F_x$, plat sur $R$, et $u_x : N \to (\varphi_*F)_y \to M_x$, qui est $R$-linéaire. Passer à la fibre en $s$, c'est réduire modulo $\mathfrak m$ : on lit l'injectivité de la clause~c), sous sa forme ponctuelle, comme l'injectivité de $N/\mathfrak m N \to \prod_x M_x/\mathfrak m M_x$, c'est-à-dire qu'un élément de $N$ dont toutes les images $u_x(n)$ sont dans $\mathfrak m M_x$ est dans $\mathfrak m N$.

\begin{encadre}
\begin{theoreme}[D'après la page 3, forme locale\piste{49-flatness-lemma-pushforward}]\label{thm:49-local}
Soient $R \to B$ un homomorphisme local d'anneaux locaux noethériens, $\mathfrak m$ l'idéal maximal de $R$, $N$ un $B$-module de type fini, $(M_x)_{x \in I}$ une famille quelconque de $R$-modules plats et $u_x : N \to M_x$ des applications $R$-linéaires. On suppose que tout $n \in N$ tel que $u_x(n) \in \mathfrak m M_x$ pour tout $x$ appartient à $\mathfrak m N$. Alors $N$ est plat sur $R$, et $(u_x)_x : N \to \prod_x M_x$ est injective. En particulier, si $u : N \to M$ est une application $R$-linéaire vers un $R$-module plat telle que $N/\mathfrak m N \to M/\mathfrak m M$ soit injective, $N$ est plat sur $R$ et $u$ est injective.
\end{theoreme}
\end{encadre}

\begin{proof}
Un anneau noethérien est cohérent (Stacks, étiquette 05CY), et sur un anneau cohérent un produit quelconque de modules plats est plat : c'est le théorème de Chase (05CZ). Le module $P = \prod_x M_x$ est donc plat sur $R$, et l'application $u = (u_x)_x : N \to P$ est $R$-linéaire. L'application $N/\mathfrak m N \to P/\mathfrak m P$ qu'elle induit est injective : si $u(n) \in \mathfrak m P$, chaque projection $u_x(n)$ est dans $\mathfrak m M_x$, puisque la projection $P \to M_x$ envoie $\mathfrak m P$ dans $\mathfrak m M_x$, et l'hypothèse donne $n \in \mathfrak m N$.

On applique alors le critère local de platitude sous la forme de l'étiquette 00ME : pour $R \to B$ local entre anneaux locaux noethériens, $N$ un $B$-module de type fini, $P$ un $R$-module plat et $u : N \to P$ une application $R$-linéaire dont la réduction modulo $\mathfrak m$ est injective, $u$ est injective et $P/u(N)$ est plat sur $R$. Le critère ne demande pas que $P$ soit de type fini sur $B$, ni que $u$ soit $B$-linéaire. On a donc une suite exacte $0 \to N \to P \to P/u(N) \to 0$ dont les deux derniers termes sont plats sur $R$ ; le premier l'est aussi (00HM), par la suite exacte des $\mathrm{Tor}$ : $\mathrm{Tor}_2^R(P/u(N), Q) \to \mathrm{Tor}_1^R(N, Q) \to \mathrm{Tor}_1^R(P, Q)$ a ses termes extrêmes nuls pour tout $R$-module $Q$. Le cas d'une seule cible est celui d'une famille à un élément.
\end{proof}

\noindent Chase, le critère local et la platitude d'un sous-module à quotient plat viennent de la bibliothèque Stacks d'AxiomMath \cite{axiomstacks}. Voir \texttt{flat\_of\_fibre\_injective\_pi} et \texttt{flat\_of\_fibre\_injective} dans \href{\leanblobmain/AlgebraicGeometry/Folder49Flatness.lean}{\texttt{Folder49Flatness.\allowbreak{}lean}}.

\begin{remarque}
La finitude de $\varphi$ aux points au-dessus de $y$ est superflue : la cible $\prod_x F_x$ n'est pas de type fini sur $\mathcal O_{Y,y}$, et la famille des points $x$ peut être infinie. Seule sert la finitude de la source : la réduction modulo $\mathfrak m$ ne détecte l'injectivité que par Nakayama, sur un module de type fini. Cette finitude-là est nécessaire : pour $R$ un anneau de valuation discrète de corps des fractions $K$, le $R$-module $K/R$ vérifie $\mathfrak m(K/R) = K/R$, de sorte que l'application nulle vers $M = 0$ satisfait l'hypothèse, et $K/R$, de torsion, n'est pas plat (remarque de l'édition, non formalisée). La platitude des $M_x$ l'est aussi, comme le montre $N = M = R/\mathfrak m$ avec $u = \mathrm{id}$. Les $u_x$ n'ont pas à être $B$-linéaires.
\end{remarque}

\begin{remarque}
Le cas d'une seule cible couvre exactement la ligne qui précède le lemme sur la page. Pour $N = \mathcal O_{A^\vee, x'}$ et $M = \mathcal O_{A, x}$, plat sur $\mathcal O_{S,s}$ parce que $A$ est plat sur $S$, les réductions modulo $\mathfrak m$ sont les anneaux locaux des fibres, $\mathcal O_{A^\vee_0, x'}$ et $\mathcal O_{A_0, x}$ ; l'homomorphisme injectif que la page écrit entre eux donne donc, si c'est la réduction d'un homomorphisme $\mathcal O_{A^\vee, x'} \to \mathcal O_{A, x}$ (celui que $\varphi_L$ induit), la platitude de $A^\vee$ sur $S$ en $x'$, sans passer par le faisceau $\varphi_*F$.
\end{remarque}

\begin{remarque}
La lecture dit que la clause~c) retenue et la ligne barrée \q{disent la même chose, la seconde ponctuellement}. Cela vaut lorsque $\varphi$ est fermé, par exemple propre. Un germe de $(\varphi_*F)_y$ dont toutes les images dans les $F_x$, $x \in \varphi^{-1}(y)$, sont nulles est une section nulle sur un ouvert $U$ contenant $\varphi^{-1}(y)$ ; si $\varphi$ est fermé, $V = Y \setminus \varphi(X \setminus U)$ est un voisinage ouvert de $y$ avec $\varphi^{-1}(V) \subset U$, et le germe est nul. Sinon, $(\varphi_*F)_y \to \prod_x F_x$ peut ne pas être injectif (pour l'immersion ouverte $\mathbf A^1 \setminus \{0\} \to \mathbf A^1$ et $y = 0$, le produit est vide et $(\varphi_*\mathcal O_X)_0$ ne l'est pas), et l'injectivité de $G_s \to \varphi_{s*}(F_s)$ au-dessus de $y$, celle de la ligne barrée, est plus faible que la forme ponctuelle démontrée ici. Dans l'application de la lecture, $\varphi_L$ est une isogénie, finie, et les deux formes coïncident.
\end{remarque}

\begin{sloppypar}
\begingroup\raggedright
\noindent\emph{Portée de la vérification.} Lean démontre le théorème~\ref{thm:49-local} sur des modules, dans \href{\leanblobmain/AlgebraicGeometry/Folder49Flatness.lean}{\texttt{Folder49Flatness.\allowbreak{}lean}}. Le passage des faisceaux aux germes n'est pas formalisé : le germe $G_y$, les applications $G_y \to (\varphi_*F)_y \to F_x$ et l'homomorphisme local $\mathcal O_{S,s} \to \mathcal O_{Y,y}$ sont des données de l'énoncé, et non construits à partir d'un morphisme de schémas. La forme barrée de la clause~c), avec l'injectivité dans $(\varphi_{s*}F_s)_y$, n'est pas couverte quand $\varphi$ n'est pas propre ; on n'a pas décidé si le lemme vaut sous cette forme.
\end{sloppypar}\par\endgroup

\begin{encadre}
\begin{theoreme*}[D'après la page 3, forme reconstituée par l'édition, critère de platitude\piste{49-flatness-lemma-pushforward}]\nonverifie
Soient $S$ localement noethérien, $\varphi : X \to Y$ un morphisme de $S$-schémas de type fini, $F$ quasi-cohérent sur $X$ et plat sur $S$, $G$ cohérent sur $Y$ et $G \to \varphi_*F$. Si $G \to \varphi_*F$ est injectif sur la fibre en $s$ aux points au-dessus de $y$, $G$ est plat sur $S$ en $y$.
\end{theoreme*}
\end{encadre}
{\small L'édition retient la forme barrée de la clause~c) ; la clause conservée, que la révision du 28 novembre 2026 lit en entier mais avec doute (\q{s'annule est nulle}), en est la forme ponctuelle. Sous la forme ponctuelle, l'énoncé se ramène au théorème~\ref{thm:49-local} une fois les germes pris ; ce passage n'est pas vérifié en Lean, et la forme barrée ne s'y ramène que si $\varphi$ est propre. On se reportera aux notes de Grothendieck (\lot{49}{1}{3}).\par}

Ne sont pas démontrés ici : la lissité de $A^\vee$ sur $S$ et l'encadrement $\dim A_0 \leqslant \dim A^\vee_0 \leqslant \dim H^1(A_0, \mathcal O_{A_0}) \leqslant \dim A_0$ des feuillets~2 et~3 ; l'isomorphisme $\Lambda^*(R^1f_*\mathcal O_A) \simeq R^*f_*\mathcal O_A$ et la liberté de $R^1f_*\mathcal O_A$ (feuillets~4 et~5) ; l'accouplement $\mathrm{Pic} \times A \to A^\vee$ et la description de son noyau (feuillets~6 à~9).
